\section{Evaluation protocols}
\label{sec:protocols}

The protocols, models, grids, metrics and statistical tests were written into a protocol file before any model was trained on the new datasets and before any test label was read. The file is self-registered, not deposited with a registry: its SHA-256 is stored with every read of test labels. The original sample was analysed in the first version with the same random split seeds, so the results on it are not confirmatory. Four deviations, all made before the first test access and all driven by validation behaviour or compute, are listed with their reasons in Appendix~\ref{app:deviations}. Test labels are returned by a single function that logs every call with the hash of the protocol file; no test label was read until every model of every protocol had been fitted (Appendix~\ref{app:access}). Each protocol is run on the full trace and on the original sample.

\begin{description}
\item[Random.] Stratified \V{prot.frac} splits into training, validation and test sets, with \V{prot.random.nsplits} split seeds. This is the protocol of the first version.
\item[Time-ordered.] Tasks are sorted by the start time of their first instance. The earliest \V{prot.time.train} train, the next \V{prot.time.val} validate and the last \V{prot.time.test} test, which on the full trace covers the last \V{full.split.time.t.testdays} days of the telemetry span. A test task is \emph{known-user} if its user has at least one training task and \emph{new-user} otherwise; we report all test tasks and each subset.
\item[User-grouped.] Five-fold \texttt{StratifiedGroupKFold}~\cite{pedregosa2011sklearn} with the user as the group (seed~0). Within the remaining users of each fold, a second grouped split selects the validation users, so no user appears in two of the three sets. We report the mean and standard deviation over the \V{prot.folds} folds and every fold.
\item[Leave-dominant-user-out.] Users are numbered by task count: User~1 has the most tasks, and User~4, fourth by tasks, is second by failures in the original sample. We remove every task of User~1, and separately of Users~1 and~4, and apply the random split with the canonical seed. This asks how much of the random-split result depends on the two users who hold most failures.
\end{description}

\paragraph{Selection.} For each split, each model's hyperparameters are chosen on the validation split, using the first initialisation seed, by validation ROC-AUC; ties within \V{stat.tie} go to the more regularised setting. The chosen setting is then used for every seed. No choice uses a test split.

\paragraph{Metrics.} ROC-AUC is the primary metric. We also report the area under the precision-recall curve (PR-AUC, computed as average precision), the Brier score~\cite{brier1950verification} and the expected calibration error (ECE) over \V{stat.bins} equal-width bins~\cite{guo2017calibration}, with reliability diagrams on the same bins. No decision threshold is tuned. A model's score on a protocol is the mean over its test sets of the mean over seeds.

\paragraph{Statistics.} Differences between \kaam{} and each other model use a paired bootstrap~\cite{efron1993bootstrap} with \V{stat.nboot} resamples. Tasks are resampled within each test set; for the user-grouped protocol whole users are resampled, since the user is the unit of generalisation there. Every model is scored on the same resamples. We report 95\% percentile intervals and two-sided bootstrap $p$-values, adjusted with Holm's method~\cite{holm1979simple} within each protocol's family of comparisons with \kaam{} (three on the full trace, where the EBMs are absent, and five on the original sample). The intervals resample test tasks with the fitted models held fixed, so they describe uncertainty about the test population and not training or fold variation; for the protocols with several test sets we also report the spread over splits or folds. A difference is \emph{not significant} when the adjusted $p$-value is at least \V{stat.alpha}. \kaam{} \emph{matches} a model when the interval of the difference lies within $\pm$\V{stat.margin} AUC. For the user-grouped protocol we also test each model's fold-mean AUC against chance with the user-level bootstrap, with Holm adjustment over the six models.
